\chapter{Estado da Arte}

\section{Árvore \textit{k}-d Dinâmicas}

A árvore k-d tree foi inicialmente introduzida pelo \citet{bentleyMultidimensionalBinarySearch1975} no artigo \textit{"Multidimensional Binary Search Trees Used for Associative Searching"} onde mostra a sua estrutura e algoritmos que podem ser utilizados, no artigo \textit{"An Algorithm for Finding Best Matches in Logarithmic Expected Time"} \citep{friedmanAlgorithmFindingBest1977} é mostrado melhorias na estrutura para buscas em tempo logaritimico, em \textit{machine learning} é utilizado algoritmo do $k$ vizinhos mais próximos (KNN) onde é utilizado e problemas de classificação e regressão, além de ser utilizados em Esembles como o KNORA e Inca-Des \citep{koDynamicClassifierSelection2008,limabarbozaIncaDesIncrementalAdaptive2024}. Essa estrutura apesar de ser boa para encontrar os k vizinhos mais próximos ela é pessima nas operações de inserção e remoção por ser uma estrutura estática, tendo isso em mente diversas abordagens foram desenvolvidas, como a árvore \textit{ikd-tree}, proposta no artigo \textit{"Ikd-Tree: An Incremental K-D Tree for Robotic Applications"} \citep{caiIkdTreeIncrementalKD2021}.

